\section{Material and Methods}
% Stub — keep original equations in a later prose PR (see PLAN.md keep-list).
% Architecture: RGB → relative depth → enhance → cues → entropy fusion → localize/merge
% → curiosity + diversity → Priority Buffer.

\subsection{System architecture}
ARTPS processes an RGB frame through relative-depth estimation, optional enhancement,
multi-cue anomaly extraction, entropy-weighted fusion, localization with tightened merge,
and ranked selection with a Priority Buffer.

\subsection{Depth-guided enhancement}
% KEEP equation: depth-guided enhancement.
% Real-ESRGAN: one paragraph only; no Results claim (unsupported without ablation).
Depth-guided enhancement modulates local contrast using the relative-depth map.
Optional Real-ESRGAN may appear as a profile-aware enhance path; it is not an ablation claim.

\subsection{Anomaly cues and entropy-weighted fusion}
% KEEP: entropy-weighted fusion.
Image, depth, and optional learned anomaly maps are combined with entropy-weighted fusion
that adapts component weights at runtime~\citep{defard2020padim,roth2022patchcore}.

\subsection{Localization, false-positive gates, and size--distance policy}
% KEEP: size--distance policy; tight merge; rover/shadow FP; object-in-shadow gate.
Localization uses tightened box merge.
Rover-body and boundary-shadow masks suppress common false positives; an object-in-shadow
gate protects shadowed rock-like support.
Size--distance policy uses image-relative relative-far and apparent size (not metric size).

\subsection{Curiosity score, diversity, and Priority Buffer}
% KEEP: curiosity score; soft similarity; uncertainty/disagreement; Priority Buffer rule.
Candidates are ranked by a curiosity score with a soft similarity penalty.
The Priority Buffer retains high raw-anomaly candidates down-weighted by diversity policy
for second-chance review when resources allow.
